\section{Each Element in Its Setting}
\label{sec:each-element}

Seven of the ten elements are Scripture and three are the Church's own
orations. The Introit joins a prayer from the Book of Daniel to the first verse
of Ps~118; the Gradual, Alleluia, Offertory and Communion each take one or two
verses from a psalm, the Communion returning to Ps~118; the lessons are St
Paul's letter to the Ephesians and St John's Gospel. Each element is set out
below in its own context, with what its words say and what the liturgical
commentators who expound it say of it, before the three readings of the whole.

\subsection{Introit: the confession of the furnace, and the blessed way}

The antiphon comes from the third chapter of Daniel. Nabuchodonosor has cast
the three young men into the furnace for refusing to worship his golden statue,
``And they walked in the midst of the flame, praising God, and blessing the
Lord'' (Dan 3:24). Then Azarias, standing in the fire, prays (3:25--45). He
blesses the God of the fathers and declares him just: ``thou art just in all
that thou hast done to us, and all thy works are true, and thy ways right, and
all thy judgments true'' (3:27). He confesses for the people, ``we have sinned,
and committed iniquity'' (3:29), and asks God not to deliver them up for ever,
for his name's sake and for the sake of Abraham, Isaac and Israel (3:34--35).
He names what exile has taken from them, ``Neither is there at this time
prince, or leader, or prophet, or holocaust, or sacrifice, or oblation, or
incense, or place of first fruits before thee'' (3:38), and offers in their
place a contrite heart: ``nevertheless, in a contrite heart and humble spirit
let us be accepted \dots\ so let our sacrifice be made in thy sight this day,
that it may please thee'' (3:39--40). Then the angel of the Lord goes down into
the furnace and drives out the flame (3:49), and the three sing the long
canticle of all creation (3:51--90). St Jerome notes in his commentary that the
Hebrew text breaks off at the youths' fall into the fire (3:23) and that the
prayer and canticle are not in the Hebrew, and he records that the Churches
read Daniel in the version of Theodotion (\work{Commentarii in Danielem}, on
3:23, and the preface).

The Missal's antiphon takes from this prayer its confession of God's just
judgment, of sin and of disobeyed commandments, and its plea for glory and
mercy, and sets them in a new order (see the note beneath the text). The prayer
it is cut from is said in part at every Mass. When the priest has offered the
chalice he bows and says Azarias's words, \latin{In spiritu humilitatis et in
animo contrito suscipiamur a te, Domine: et sic fiat sacrificium nostrum in
conspectu tuo hodie, ut placeat tibi, Domine Deus} (1962 \work{Ordo Missae},
no.~1033).

The verse is the first of Ps~118, whose twenty-two strophes of eight verses
follow the letters of the Hebrew alphabet and praise the law of God. Its
first word is a beatitude. The verse sets the way, \latin{in
via}, beside the law, \latin{in lege Domini}, and calls blessed those who walk
in it undefiled.

The medieval liturgical commentators Sicard of Cremona (\work{Mitrale}
VIII.20) and William Durandus (\work{Rationale} VI.137) hear in the Introit
the voice of Daniel recalling the captivity and ascribing it to the people's
deserts by God's just judgment. Bl.\ Ildefonso Schuster, the Benedictine
liturgist and Cardinal Archbishop of Milan, observes that the verses are from
Daniel ``but the actual words of the text are not quoted'', and he reads the
antiphon as ``the sorrowful confession of guilt which leads back the sinner to
the way of reconciliation''; in its last clause he finds ``the hope which
underlies the act of contrition'' (\work{The Sacramentary} III, pp.~174--175).

\subsection{Collect: pardon and peace, for a mind at rest}

The Collect asks two gifts and names two effects. \latin{Indulgentiam} and
\latin{pacem}, pardon and peace, are asked of God \latin{placatus}, as one
appeased; and they are asked so that the faithful may be cleansed from every
offence, \latin{ab omnibus mundentur offensis}, and serve God \latin{secura
mente}, with a mind free of care. The order is pardon first, then peace; and
the service it ends with is the service of minds at rest. Schuster expounds
that order: ``Peace comes after mercy, for as long as grace has not blotted out
the sin, the heart, torn by remorse, weakened by its own passions and at war
with itself, cannot find peace.'' He sets beside it \latin{Non est pax impio},
St Paul's name for Christ, \latin{Pax nostra}, and the beatitude of the clean
of heart, and he sees the fruit of the twofold gift in ``a great ease and
facility in the performing of good actions'' (p.~175).

\subsection{Epistle: walking wisely in evil days}

The lesson comes in the middle of St Paul's call to the Ephesians to live as
light. ``For you were heretofore darkness, but now light in the Lord. Walk then
as children of the light'' (Eph 5:8); and just before the lesson, ``Rise, thou
that sleepest, and arise from the dead: and Christ shall enlighten thee''
(5:14). The verses the Missal reads draw the consequence, ``See therefore'': walk
carefully, as wise, redeeming the time, because the days are evil; do not be
unwise but understand the will of God; be filled not with wine but with the
Spirit, speaking in psalms, hymns and spiritual canticles, singing in the heart,
giving thanks always for all things, and subject one to another in the fear of
Christ (5:15--21). The last of these turns the letter to the household, to
husbands and wives as Christ and the Church, ``This is a great sacrament: but I
speak in Christ and in the church'' (5:22--33).

St Thomas Aquinas gives two expositions of ``redeeming the time''. In the first,
which he judges the better, a man buys back his time as one buys off a vexation,
by giving up something of his own right, because since Adam snares lie ready
on every side; in the second, he makes up the time lost in sin by doing more
good works now (\work{In Epistolam ad Ephesios} c.~5, lect.~6). Schuster hears
the lesson in its season: ``The season is already advanced, the vintage is now
over, it is necessary to redeem the time which has been spent unprofitably to
the soul, and instead of wasting it on feasting \dots\ we should prepare for
evil days---that is, for old age and death by adding to our good works while yet
we have time, through the grace of the Holy Ghost'' (p.~175).

\subsection{Gradual: the eyes of all, and food in due season}

Ps~144, titled \latin{Laudatio ipsi David}, is an alphabetic hymn of God's
kingship, and its second half praises his care for every creature: ``The Lord
lifteth up all that fall: and setteth up all that are cast down'' (v.~14); then
the Gradual's two verses; then ``The Lord is just in all his ways'' (v.~17) and
``The Lord is nigh unto all them that call upon him'' (v.~18). The Gradual
sings of food given to all, at the right time, from an open hand.

St John Chrysostom expounds the verses of God's providence over all things
(\work{Expositio in Psalmos}, on Ps~144:15--16). It is not rain, earth or air
that brings the fruits to ripeness but God's command, as the sun rises on the
evil and the good; ``in due season'' means either that each thing is given in
its own season of the year, so that the farmer may rest and the produce not
spoil, or that God gives food to those who need it; and the open hand is his
working power, set against those who worshipped the air and the sun as the
givers of fruit. St Hilary of Poitiers, on the same verses, says that those
who are cast down look to God with all their eyes, since every support of life
is his, both the food of the body in its season and the filling of souls with
spiritual gifts (\work{Tractatus super Psalmos}, in Ps.~144). Cassiodorus
reads \latin{in tempore opportuno} as a warning that what is asked is not always
given at once: many who ask are happily deferred, and the one who prays should
ask God to give what he knows to be expedient (\work{Expositio Psalmorum}
144.15). Schuster finds in the verse ``a marked eucharistic significance'':
``Holy Communion is the true universal bread which God prepares for his
children in every region of the earth'' (pp.~175--176).

\subsection{Alleluia: a heart that is ready}

Ps~107, titled \latin{Canticum Psalmi ipsi David}, is made of the ends of two
other psalms: its verses 2--6 repeat Ps~56:8--12, and its verses 7--14 repeat
Ps~59:7--14. The Alleluia's verse is therefore also the verse in which the
psalmist of Ps~56, whose title sets it in David's flight from Saul into the
cave and who has been delivered ``from the midst of the young lions'' (56:1,
5), turns to praise: ``My heart is ready, O God, my heart is ready: I will sing,
and rehearse a psalm. Arise, O my glory, arise psaltery and harp: I will arise
early'' (56:8--9). St Augustine gives Ps~107 no exposition of its own, since it
repeats two psalms already expounded, and reads its title as pointing to
prophecy. Cassiodorus hears Christ the King speaking throughout: the readiness
of heart is the holiness of the Incarnation, in which no discord of sin hindered
his praise of the Father, and the faithful say the verse rightly when they
trust that their sins are forgiven by his gift; the dawn at which the singer
will rise is the Resurrection (\work{Expositio Psalmorum} 107.2--3). St Robert
Bellarmine, on the same words in Ps~56, makes them a soul's whole surrender:
``ready to live, ready to die, ready to rule, ready to be trampled on, ready to
take anything cheerfully from your hand'' (\work{A Commentary on the Book of
Psalms}, Ps~56:8). Schuster reads the Missal's \latin{gloria mea} of God: ``God
is called the glory of the soul'', as author of its glory and as ``the just
judge of our merits'' (p.~176).

\subsection{Gospel: the second sign at Cana}

The healing of the ruler's son closes the fourth chapter of St John. It
follows the two days that Jesus spent among the Samaritans, of whom ``many more
believed in him, because of his own word'' and who confessed him ``indeed the
Saviour of the world'' (Jn 4:41--42). He goes on into Galilee, ``For Jesus
himself gave testimony that a prophet hath no honour in his own country''
(4:44); the Galileans receive him, ``having seen all the things he had done at
Jerusalem on the festival day'' (4:45). He comes to Cana, ``where he made the
water wine'' (4:46), and there the ruler from Capharnaum finds him. The
evangelist closes the episode, ``This is again the second miracle that Jesus
did, when he was come out of Judea into Galilee'' (4:54). The Missal's lesson
begins after the mention of Cana and ends before the count of signs, so that
it reads the healing alone.

The ruler, \latin{regulus}, is a man of rank. St John Chrysostom takes him to
be ``of royal race, or possessed some dignity from his office'', and St Thomas
Aquinas an official of Herod's household or an officer of the Emperor (\work{Super
Ioannem} c.~4, lect.~7). Both deny that he is the centurion of St Matthew's
eighth chapter; Aquinas counts four differences, in the sickness (a palsy, here
a fever), the sufferer (a servant, here a son), the petition (the centurion
begged Christ not to come, the ruler begged him to come down) and the place.
The servants' answer, ``Yesterday at the seventh hour'', marks the hour of the
word, and the evangelist's last words reach the whole household.

St Gregory the Great preached on this Gospel at the basilica of the martyrs
Nereus and Achilleus on their feast (\work{Homiliae in Evangelia} 28, heading),
and the readings below take up his homily. Among the medieval commentators,
Rupert of Deutz reads the fever that had all but extinguished the boy as the
shaking of a noble and fighting soul in temptation, as Job's was (\work{De
divinis officiis} XII.21). William Durandus makes the ruler a figure of kings
and prelates, who ought to be humble, and says that he believed the Lord could
heal his son but not that he was God, since he believed that he could not heal
unless he were present (\work{Rationale} VI.138). St Anthony of Padua, in his
Sunday sermon on this Gospel, makes every faithful person a ruler over himself
and his son the soul, which falls sick in Capharnaum, a name he interprets as a
field of fatness or of consolation.

\subsection{Offertory: by the rivers of Babylon}

Ps~136 is the song of the Babylonian exile, titled in the Vulgate
\latin{Psalmus David, Ieremiae}; St Hilary notes that among ``the Hebrew
elders'' it has no title. Its nine verses move from lament to refusal to oath
to imprecation. The captives sit by the rivers and weep, remembering Sion
(v.~1); they hang their instruments on the willows (v.~2); their captors demand
``a hymn of the songs of Sion'', and they answer, ``How shall we sing the song
of the Lord in a strange land?'' (vv.~3--4). Then one voice swears: ``If I
forget thee, O Jerusalem, let my right hand be forgotten'' (vv.~5--6). The psalm
ends by asking God to remember the children of Edom who cried ``Rase it, rase
it'', and by blessing whoever repays the daughter of Babylon and dashes her
little ones against the rock (vv.~7--9). The Offertory sings the first verse
only, in a form that speaks to Sion, \latin{tui, Sion}; the next readings take
up the rest of the psalm in the Fathers' expositions. St Robert Bellarmine
places the scene in the country around the city, where the captives were set
to labour along the many rivers of the province, ``for it is well known that
the Euphrates was the only river that ran through it''. Schuster takes the
waters as ``the passions in which the unhappy sinner quenches his thirst,
whilst the faithful disciple sits sadly beside that tainted stream'' (p.~177).

\subsection{Secret: heavenly medicine for the heart}

The Secret is a prayer for healing. It asks that \latin{haec mysteria}, these
mysteries, furnish heavenly medicine and purge the vices of the heart. Its
vocabulary meets the Gospel's: the son \latin{infirmabatur}, was sick; he
\latin{incipiebat mori}, was beginning to die; the fever, \latin{febris}, left
him. Schuster reads the prayer of the Eucharist, asking ``that the holy
Sacrament may be to us a spiritual medicine and antidote, to protect us from
the virus of sin which poisons our blood'' (p.~177). The Missal follows it with
the direction for the Preface of the Most Holy Trinity.

\subsection{Communion: remember thy word}

The Communion takes the first two verses of the seventh strophe of Ps~118,
\latin{Zain}. The psalmist asks God to remember the word in which he gave hope,
and says that this hope has comforted him in his humiliation; the verse goes on,
in the Vulgate, ``because thy word hath enlivened me'' (v.~50), and the strophe
continues, ``The proud did iniquitously altogether: but I declined not from thy
law. I remembered, O Lord, thy judgments of old: and I was comforted''
(vv.~51--52). Every witness read on this verse takes ``remember'' as a prayer
that God fulfil a promise and not as a reminder to one who could forget: St
Augustine, St Hilary, Cassiodorus, St Robert Bellarmine and Durandus. Hilary
sharpens it: the prophet does not remind God, but asks to be found worthy that
God should remember his word in him (\work{Tractatus super Psalmos}, in
Ps.~118, \latin{Zain}). Schuster, at the moment of Communion, turns the word
into the Word: ``let us offer to the eternal Father this Word, in whom he is
well-pleased, and who is the cause of all our hope'' (p.~177).

\subsection{Postcommunion: made worthy by obedience}

The Postcommunion asks that those who have received may be made worthy of the
sacred gifts, and it names the way: \latin{fac nos \dots\ tuis semper
oboedire mandatis}, make us always obey thy commandments. Its noun and verb are
those of the Introit, which confessed \latin{mandatis tuis non oboedivimus}, so
that the Mass ends by asking for the obedience it began by confessing it had not
given. The continuation of Dom Prosper Guéranger's \work{The Liturgical Year},
written by Dom Lucien Fromage and published in English in 1883, reads it so:
``A persevering fidelity in observing God's commandments is the best
preparation a christian can make for approaching to the holy Table''; Schuster
asks of it ``our habitual docility to the inspirations of the Holy Ghost''
(pp.~177--178).
